\subsubsection{Execution Algorithms}

\begin{principle} \hlt{Algo as Layered Decision System} \cite{lehalle_2018}\\
An execution algorithm is not only a schedule. It is a layered system that turns a parent order into child orders, adapts to market state, and records enough context to explain the realised cost afterwards.
\end{principle}

\begin{taxonomy} \hlt{Algorithm Layer Segmentation}
\begin{enumerate}[label=\roman*.]
\setlength{\itemsep}{0pt}
\item \hlt{execution parameters}: side, size, benchmark, timing, limits, participation, urgency, and venues;
\item \hlt{risk-control layer}: sets the curve from impact, timing risk, alpha, and context;
\item \hlt{tactical layer}: uses SORs, liquidity seeking, pegging, and dark probes to realise it.
\end{enumerate}
Risk control sets the envelope; tactics work inside it.
\end{taxonomy}

\begin{process} \hlt{Execution Decision Loop}
\begin{enumerate}[label=\roman*.]
\setlength{\itemsep}{0pt}
\item initialise the parent order and benchmark;
\item compute a target trajectory from historical profiles and model parameters;
\item observe quotes, trades, fills, rejects, cancellations, and venue responses;
\item revise child-order size, price, urgency, and venue split;
\item publish monitoring and performance indicators for the trader.
\end{enumerate}
\end{process}

\begin{taxonomy} \hlt{Pre-Trade Segmentation}
\begin{enumerate}[label=\roman*.]
\setlength{\itemsep}{0pt}
\item \hlt{expected cost}: spread, impact, volatility, and typical quantity on the book;
\item \hlt{profile}: expected volume curve, duration, and participation bounds;
\item \hlt{context}: news, volatility regime, sector, country, currency, and event risk;
\item \hlt{portfolio coupling}: common end times, exposure constraints, and cross-order synchronisation.
\end{enumerate}
\end{taxonomy}

\begin{table}[htbp]
\centering
\scriptsize
\stepcounter{table}
\renewcommand{\arraystretch}{1.14}
\setlength{\tabcolsep}{4pt}
\begin{tabularx}{\textwidth}{@{}>{\RaggedRight\arraybackslash}p{0.14\textwidth} >{\RaggedRight\arraybackslash}p{0.17\textwidth} >{\RaggedRight\arraybackslash}p{0.31\textwidth} >{\RaggedRight\arraybackslash}X@{}}
\hline
\textbf{Benchmark} & \textbf{Stock type} & \textbf{Trade type} & \textbf{Main feature} \\
\hline
PoV & Medium to large market depth & Long-duration position trade & Flow-following; reactive within participation bounds \\
VWAP/TWAP & Any market depth & Hedging order; long-duration position; tracking-error unwind & Volume-curve schedule; passive when possible \\
IS (Implementation Shortfall) & Medium liquidity depth & Alpha extraction; non-linear hedge, typically gamma; inventory-driven trade & Fast on opportunity; cuts losses on adverse moves \\
Liquidity seeker & Poor fragmented market depth & Alpha extraction; opportunistic position mounting; already split or scheduled order & Pool search; dispersed liquidity; low signalling \\
\hline
\end{tabularx}

\vspace{0.25em}
{\footnotesize Table \thetable. Typical features of trading algorithms.}
\end{table}

\begin{practitionernote} \hlt{What-If Controls}\\
Useful tools support live counterfactuals: finish now, slow down, switch benchmark, seek dark liquidity, or cap participation. The aim is disciplined comparison before changing trajectory.
\end{practitionernote}
